\part{Bugs, and how they were found}

\noindent\textbf{\large What this part covers}

\medskip
\noindent
Real failures from building this system. Every one of them actually happened.
Each is told in the same four steps: the \textbf{symptom}, the \textbf{wrong
guess} that came first, the \textbf{evidence} that corrected it, and the
\textbf{fix}.

\medskip
\noindent
The wrong guesses are included deliberately. Debugging is taught as though the
right explanation arrives immediately; it does not. What separates people who
are good at it is not better hunches --- it is refusing to act on a hunch before
checking it.

\chapter{Bugs in the application}
\label{ch:bugs-frontend}

\section{The button that broke every test}

\textbf{Symptom.} A reveal toggle was added to the password field. All seven
browser tests failed instantly, with an error about \emph{strict mode}.

\textbf{Wrong guess.} The tests are stale --- the page changed, so they need
updating. Reasonable, and wrong in an important way: it treats the failure as
noise instead of reading it.

\textbf{Evidence.} Playwright's message named the problem exactly: the locator
resolved to \emph{two} elements.

\where{web/e2e/incident-lifecycle.spec.ts}
\begin{lstlisting}[language=tsx]
await page.getByLabel('Password', { exact: true }).fill(DEMO_PASSWORD);
// Exact, because the field shares the word "password" with the reveal
// toggle's aria-label ("Show password"). getByLabel substring-matches by
// default, so without this the locator resolves to two elements and trips
// strict mode -- which is exactly what happened when the toggle was added.
\end{lstlisting}

\textbf{The cause.} \fpath{getByLabel('Password')} matches \emph{substrings}.
The new button's accessible name is ``Show password'', which contains
``password''. Two elements, one locator.

\textbf{The fix.} \fpath{\{ exact: true \}}.

\note{Strict mode is Playwright refusing to guess. A less careful tool would
have picked the first match and passed --- and the test would silently have been
typing the password into a button from then on. An error that stops you is
better than a pass that means nothing, and this is worth remembering when a tool
feels pedantic.}

\section{The manual that would have lied}

\textbf{Symptom.} The script that captures the screenshots for this book
reported success --- 21 files written. Four of them were byte-identical to the
login page, including the one labelled ``the metrics page''.

\textbf{Wrong guess.} A timing problem: the screenshot fired before the page
finished rendering. Plausible, and it produced a fix that did not work.

\textbf{Evidence.} Two things, in order. First, file sizes: four files at
\emph{exactly} 464~KB, the same as the login screenshot. Identical sizes mean
identical images. Second, a trace of the actual network calls, which showed the
session being lost partway through the run.

\textbf{The cause.} Two facts that are harmless separately:

\begin{itemize}[leftmargin=1.4cm]
\item The access token is held in memory only --- so a full page reload loses it
  and the app must call \fpath{/auth/refresh} to get another.
\item \fpath{/auth/refresh} is rate-limited against the \emph{same} per-identity
  bucket as \fpath{login}.
\end{itemize}

The capture script navigated with \fpath{page.goto()}, which is a full reload, so
every screen cost one refresh. Add a deliberate failed login for the
error-message screenshot --- using the same account --- and the tour ran out of
budget partway through, was logged out, and cheerfully photographed the login
page over and over.

\textbf{The fix.} Two changes. Navigate by clicking the sidebar, as a user does,
which needs no reloads at all. And spend the deliberate failure on a
\emph{different} account.

But the important change was neither of those:

\where{web/capture-screens.mjs}
\begin{lstlisting}[language=tsx]
if (expect) {
  await page.getByRole('heading', { name: expect }).first()
    .waitFor({ state: 'visible', timeout: 20000 })
    .catch(() => {
      throw new Error(
        `REFUSING to save ${name}.png -- expected heading "${expect}" never appeared. ` +
          `The page is probably showing Sign in instead.`,
      );
    });
}
\end{lstlisting}

\warnbox{The script now refuses to save a screenshot it cannot verify. This is
the general lesson: a process that produces output nobody checks will eventually
produce confident, wrong output. A picture is worse than prose here, because a
reader trusts a screenshot completely --- they assume a photograph cannot be
mistaken about what it shows.}

\chapter{Bugs in the backend}
\label{ch:bugs-backend}

\section{The one that only failed on PostgreSQL}

\textbf{Symptom.} All tests passed locally. The same commit failed in CI with
\fpath{SQLSTATE[25P02]}: \emph{current transaction is aborted, commands ignored
until end of transaction block}.

\textbf{Wrong guess.} A flaky test, or a CI configuration difference. It was
neither --- it was a real bug that local testing structurally could not find.

\textbf{Evidence.} The tests run against in-memory SQLite for speed; CI also
runs them against PostgreSQL. Only PostgreSQL has this behaviour.

\textbf{The cause.} The API prevents duplicate submissions with an
\emph{idempotency key}: it inserts a row, and relies on a unique index to reject
a second attempt. Catching that rejection is the mechanism.

But in PostgreSQL, \textbf{a failed statement poisons the whole transaction}.
After the duplicate-key error, every subsequent statement fails with 25P02 until
someone rolls back. SQLite does not do this, so locally the code worked
perfectly.

The real-world effect: a user who double-clicked ``Report incident'' got
\fpath{500 Internal Server Error} rather than the polite ``you already submitted
this'' the feature was written to produce.

\textbf{The fix.} Wrap the insert in a nested transaction, which PostgreSQL
implements as a \emph{savepoint}. The failure then rolls back only to the
savepoint, and the outer transaction survives.

\where{api/app/Http/Middleware/HandleIdempotency.php}
\begin{lstlisting}[language=PHP]
return DB::transaction(fn (): IdempotencyKey => IdempotencyKey::query()->create([...]));
\end{lstlisting}

\note{One line, and it is not obvious from reading it that it is doing anything.
Laravel turns a nested \fpath{DB::transaction} into \fpath{SAVEPOINT}, so the
error is contained. Without the surrounding story, this line looks like it could
be deleted.}

\warnbox{The general lesson is the one worth carrying: \textbf{a fast local test
database that is not the production database will hide a class of bugs, and it
will hide precisely the ones that involve transactional behaviour.} The
substitute is chosen for speed, and speed usually means fewer guarantees ---
which is the same thing as fewer failure modes to reproduce.}

\section{The test that proved nothing}

\textbf{Symptom.} A newly-added test passed locally and failed in CI ---
\emph{different} CI job from the last story, and a different cause.

\textbf{Wrong guess.} The rate limit is too low on CI. Checked the logs for
\fpath{429} responses: none. The guess was wrong, and checking took two minutes.

\textbf{Evidence.} The environment blocks differ:

\begin{lstlisting}
phpunit.xml           CACHE_STORE = array     fresh for every test
CI, PostgreSQL job    CACHE_STORE = redis     shared for the whole run
\end{lstlisting}

The rate limiter keeps its counters in the cache. Nine authentication requests
happen in that test file \emph{before} the new test runs. With a per-test cache
that does not matter. With a shared one, the per-IP budget was nearly spent and
the first attempt in the new test came back 429 --- failing an assertion about a
different limit entirely.

\textbf{The fix.} Reset the limiter at the start of the test, so it counts from
a known state on any driver. Then --- and this is the part that mattered ---
\textbf{prove the fix}: run the whole suite against a real Redis, with and
without the change, flushing between runs.

\begin{lstlisting}
with the reset      OK (114 tests, 547 assertions)
without the reset   FAILURES! 1  -- the exact CI failure
\end{lstlisting}

\note{Notice the shape repeating from the previous story. Local uses a
\emph{more forgiving substitute} --- in-memory SQLite, in-memory cache --- and
both forgive exactly the thing that breaks elsewhere. Test doubles are chosen for
speed, and speed almost always means \emph{isolation}; isolation is precisely the
property the real component lacks.}

\section{The header that was documented and never sent}

\textbf{Symptom.} None. Nothing failed. This was found by writing a test for
behaviour nobody had tested.

\textbf{Evidence.} The API's own error message for rate limiting read:
\emph{``Too many requests. Slow down and retry after the period in the
Retry-After header.''} The header was not being sent.

\textbf{The cause.} The central error renderer rebuilt every response from the
status code and body alone --- discarding any headers the exception carried. The
throttler attaches \fpath{Retry-After} and the \fpath{X-RateLimit-*} trio to its
exception, and all of them were silently dropped.

\textbf{The fix.}

\where{api/app/Exceptions/ApiExceptionRenderer.php}
\begin{lstlisting}[language=PHP]
/*
 * Carry over headers the exception was throwing *with*, not just its
 * status. The throttler is the case that matters: it raises a 429
 * carrying Retry-After and the X-RateLimit-* trio, and rebuilding the
 * response from scratch silently dropped all of them.
 */
$headers = $e instanceof HttpExceptionInterface ? $e->getHeaders() : [];

return new JsonResponse($payload, $status, $headers);
\end{lstlisting}

\warnbox{There was no rate-limiting test at all. The limiters were defined and
attached, which \emph{reads} as correct --- and nothing proved they refused
anything. This is the most common kind of hidden defect in a codebase: not code
that is wrong, but code that is never executed by anything that would notice.}

\section{The failure that turned a success into a 500}

\textbf{Symptom.} During a CI run with a misconfigured Redis, creating an
incident returned \fpath{500}.

\textbf{The cause.} The incident was written and committed successfully. Then the
code tried to queue the notification job, Redis was unreachable, the exception
propagated, and the request returned 500.

So the incident \emph{existed} --- and the reporter was told it had failed. On a
\fpath{sev1} at 3am, that person now re-reports it.

\textbf{The fix.} Catch it, log it, and leave the row \fpath{pending}.

\where{api/app/Services/Notifications/NotificationDispatcher.php}
\begin{lstlisting}[language=PHP]
try {
    SendIncidentNotification::dispatch($notification->getKey());
} catch (Throwable $e) {
    /*
     * The queue broker is unreachable. This must not propagate.
     *
     * By the time we get here the incident and its notification rows
     * are already committed -- dispatch() runs after the transaction
     * on purpose. Letting this throw would return 500 for a write
     * that actually succeeded, telling a reporter their SEV-1 was
     * not filed when it was.
     */
    Log::error('notification.enqueue_failed', [...]);
}
\end{lstlisting}

\note{A scheduled command already existed to sweep notifications stuck in
\fpath{pending}. The safety net had been built --- it was simply unreachable,
because the failure it was designed to catch killed the request before the row
could be left in the state it looks for. Building a recovery mechanism and then
failing loudly instead of entering the state it recovers from is a surprisingly
common mistake.}
